\chapter{Acknowledgments}

First and foremost, I'd like to thank my advisor, Aise Johan de Jong, for his support and guidance throughout my PhD. 
Speaking with him always led to an incredibly stimulating intellectual discussion, and I am very grateful for the many hours he spent helping me develop my ideas. 
I came into this PhD with a preconceived notion of the mathematics that I liked and disliked, 
and have Johan to thank for helping me break out of that mindset and develop an appreciation for many areas of mathematics,
 because one never knows when a particular technique from a seemingly unrelated area of mathematics will be useful for solving a problem.
Johan also taught me the importance of being precise and careful with my mathematical writing, a skill which I have also 
found useful beyond mathematics. With no doubt, I will be drawing on these skills for the rest of my life, and I am very grateful to Johan for helping me develop them.\\
\indent I would also like to thank my former advisor, James Borger. Even after my Master's thesis, James continued to be a great source of support and encouragement, 
and my time under his supervision shaped my mathematical interests and approach to mathematics in many ways. Most of the material on semirings in this thesis
was inspired by collaboration with James, and I am still in awe of James's insights and creativity in developing beautiful mathematical theory.\\
\indent Next, I would like to thank the whole Columbia Math Department support staff for their incredible work in keeping the department running smoothly. 
In particular, I would like to thank Nathan Schweer for his help with many administrative issues I faced during my PhD, as well as his support and encouragement.\\
\indent I would also like to thank Juan Esteban Rodríguez Camargo and Hanlin Cai for many helpful discussions. 
Juan was very generous in helping me understand the literature on $\infty$-categories, descendability, and with proving Proposition~\ref{proposition:symformula}. 
Hanlin was very helpful in discussing ideas concerning Chapter~\ref{chapter:affine}. In particular, Hanlin frequently suggested
that considering the absolute integral closure of a Noetherian local ring may be useful, and he was the first to rephrase the equicharacteristic 
proof of Theorem~\ref{thm:affine-mixed} in terms of the Riemann-Hilbert functor of Bhatt and Lurie. One of the many threads of our discussion 
was the `potential strategy' to prove Theorem~\ref{thm:affine-mixed} by tilting to equicharacteristic outlined in the introduction of Chapter~\ref{chapter:affine}. 
It would've been incredibly satisfying if these arguments had worked, but unfortunately we have not been able to do so.\\
\indent Last, but certainly not least, I'd like to thank my fiancée Liberty for her love and support throughout my whole PhD. Life in New York City
is both wonderful and challenging, and I am very grateful to have had Liberty by my side to share in the wonderful moments and to support me through the challenging ones.\\
\indent My family, especially my parents, have also been a great source of support and encouragement throughout my life, and I couldn't have 
made it this far without them. Thank you to my sister Eve for coming up with the name `indivisible sequences'.\\
\indent Thank you to the Fulbright program for providing me with the funding to pursue my PhD at Columbia University.